%! Equivalent Representations of Trigonometric Functions — Unit 7, Lesson 7.5 — Warm-Up (Answer Key)
\documentclass[10pt]{article}
\usepackage{precalculus-article}
\usepackage{precalculus-key}   % pulls in -boxes; do NOT also load -boxes
\newcommand{\TallMath}[1]{$\displaystyle #1\rule[-1.4em]{0pt}{3.2em}$}

\begin{document}

\pageheader{Unit 7, Lesson 7.5}{Warm-Up --- Answer Key}
\namedateperiod

\vspace{0.10in}

\begin{notesbox}{Spiral Review --- Prerequisites for Lesson 7.5}

\textbf{1.}\quad Evaluate $\sin^2\!\left(\dfrac{\pi}{6}\right) + \cos^2\!\left(\dfrac{\pi}{6}\right)$
using exact unit-circle values.

\vspace{4pt}
$\sin\!\left(\tfrac{\pi}{6}\right) = $ \ans{$\dfrac{1}{2}$}\quad
$\cos\!\left(\tfrac{\pi}{6}\right) = $ \ans{$\dfrac{\sqrt{3}}{2}$}

\vspace{6pt}
$\sin^2\!\left(\tfrac{\pi}{6}\right) + \cos^2\!\left(\tfrac{\pi}{6}\right)
= $ \ans{$\dfrac{1}{4}$} $+$ \ans{$\dfrac{3}{4}$} $=$ \ans{$1$}

\vspace{0.14in}
\textbf{2.}\quad Expand $(a + b)^2$ algebraically.

\vspace{4pt}
$(a + b)^2 = $ \ans{$a^2 + 2ab + b^2$}

\vspace{0.14in}
\textbf{3.}\quad Recall exact unit-circle values.

\vspace{4pt}
$\sin\!\left(\tfrac{\pi}{6}\right) = $ \ans{$\tfrac{1}{2}$}\hspace{0.6cm}
$\cos\!\left(\tfrac{\pi}{6}\right) = $ \ans{$\tfrac{\sqrt{3}}{2}$}\hspace{0.6cm}
$\sin\!\left(\tfrac{\pi}{4}\right) = $ \ans{$\tfrac{\sqrt{2}}{2}$}\hspace{0.6cm}
$\cos\!\left(\tfrac{\pi}{4}\right) = $ \ans{$\tfrac{\sqrt{2}}{2}$}

\vspace{0.14in}
\textbf{4.}\quad Simplify: $\dfrac{\sin\theta}{\cos\theta} = $ \ans{$\tan\theta$}

\begin{teachernote}
Q1 motivates the Pythagorean identity: the sum always equals 1 regardless of the angle.
Q3 primes students for the sum-identity example $\sin(45^\circ + 30^\circ)$.
Q4 primes students for the identity $\tan\theta \cdot \cos\theta = \sin\theta$.
\end{teachernote}

\end{notesbox}

\end{document}
